% Options for packages loaded elsewhere
\PassOptionsToPackage{unicode}{hyperref}
\PassOptionsToPackage{hyphens}{url}
\PassOptionsToPackage{dvipsnames,svgnames,x11names}{xcolor}
\documentclass[
  10pt,
]{article}
\usepackage{xcolor}
\usepackage[margin=1cm,top=1cm,bottom=1cm,left=1cm,right=1cm,includeheadfoot]{geometry}
\usepackage{amsmath,amssymb}
\setcounter{secnumdepth}{3}
\usepackage{iftex}
\ifPDFTeX
  \usepackage[T1]{fontenc}
  \usepackage[utf8]{inputenc}
  \usepackage{textcomp} % provide euro and other symbols
\else % if luatex or xetex
  \usepackage{unicode-math} % this also loads fontspec
  \defaultfontfeatures{Scale=MatchLowercase}
  \defaultfontfeatures[\rmfamily]{Ligatures=TeX,Scale=1}
\fi
\usepackage{lmodern}
\ifPDFTeX\else
  % xetex/luatex font selection
  \setmainfont[]{DejaVu Serif}
  \setmonofont[]{DejaVu Sans Mono}
\fi
% Use upquote if available, for straight quotes in verbatim environments
\IfFileExists{upquote.sty}{\usepackage{upquote}}{}
\IfFileExists{microtype.sty}{% use microtype if available
  \usepackage[]{microtype}
  \UseMicrotypeSet[protrusion]{basicmath} % disable protrusion for tt fonts
}{}
\usepackage{setspace}
\makeatletter
\@ifundefined{KOMAClassName}{% if non-KOMA class
  \IfFileExists{parskip.sty}{%
    \usepackage{parskip}
  }{% else
    \setlength{\parindent}{0pt}
    \setlength{\parskip}{6pt plus 2pt minus 1pt}}
}{% if KOMA class
  \KOMAoptions{parskip=half}}
\makeatother
\usepackage{listings}
\newcommand{\passthrough}[1]{#1}
\lstset{defaultdialect=[5.3]Lua}
\lstset{defaultdialect=[x86masm]Assembler}
\usepackage{longtable,booktabs,array}
\usepackage{caption}
\captionsetup[table]{skip=6pt}
\newcounter{none} % for unnumbered tables
\usepackage{calc} % for calculating minipage widths
% Correct order of tables after \paragraph or \subparagraph
\usepackage{etoolbox}
\makeatletter
\patchcmd\longtable{\par}{\if@noskipsec\mbox{}\fi\par}{}{}
\makeatother
\setlength{\emergencystretch}{3em} % prevent overfull lines
\providecommand{\tightlist}{%
  \setlength{\itemsep}{0pt}\setlength{\parskip}{0pt}}
% Enable graphics inclusion and ensure figure numbering works
\usepackage{graphicx}
\renewcommand{\figurename}{Figure}

% Configure fonts for Unicode support with fallbacks
\usepackage{newunicodechar}
\newunicodechar{⁴}{\textsuperscript{4}}
\newunicodechar{₄}{\textsubscript{4}}

% Math notation macros
\newcommand{\norm}[1]{\lVert #1\rVert}
\newcommand{\abs}[1]{\lvert #1\rvert}

% Enhanced code block styling for better contrast and readability
\usepackage{fancyvrb}
\usepackage{xcolor}
\usepackage{listings}

% Define custom colors for code blocks
\definecolor{codebg}{RGB}{245, 245, 245}      % Light gray background
\definecolor{codeborder}{RGB}{200, 200, 200}  % Medium gray border
\definecolor{codefg}{RGB}{50, 50, 50}         % Dark gray text

% Configure Verbatim environment for inline code
\DefineVerbatimEnvironment{Verbatim}{Verbatim}{%
    fontsize=\small,
    frame=single,
    framerule=0.5pt,
    framesep=3pt,
    rulecolor=\color{codeborder},
    bgcolor=\color{codebg},
    fgcolor=\color{codefg}
}

% Configure code block styling
\DefineVerbatimEnvironment{Highlighting}{Verbatim}{%
    fontsize=\footnotesize,
    frame=single,
    framerule=0.5pt,
    framesep=5pt,
    rulecolor=\color{codeborder},
    bgcolor=\color{codebg},
    fgcolor=\color{codefg}
}

% Style inline code with \texttt
\renewcommand{\texttt}[1]{%
    \colorbox{codebg}{\color{codefg}\ttfamily #1}%
}

% Configure listings package for code blocks
\lstset{
    backgroundcolor=\color{codebg},
    basicstyle=\footnotesize\ttfamily\color{codefg},
    breakatwhitespace=false,
    breaklines=true,
    captionpos=b,
    commentstyle=\color{codefg},
    deletekeywords={...},
    escapeinside={\%*}{*)},
    extendedchars=true,
    frame=single,
    framerule=0.5pt,
    framesep=5pt,
    keepspaces=true,
    keywordstyle=\color{codefg},
    language=Python,
    morekeywords={*,...},
    numbers=left,
    numbersep=5pt,
    numberstyle=\tiny\color{codefg},
    rulecolor=\color{codeborder},
    showspaces=false,
    showstringspaces=false,
    showtabs=false,
    stepnumber=1,
    stringstyle=\color{codefg},
    tabsize=2,
    title=\lstname
}

% Override any Pandoc default lstset configurations
\AtBeginDocument{
    \lstset{
        backgroundcolor=\color{codebg},
        basicstyle=\footnotesize\ttfamily\color{codefg},
        frame=single,
        framerule=0.5pt,
        framesep=5pt,
        rulecolor=\color{codeborder},
        numbers=left,
        numbersep=5pt,
        numberstyle=\tiny\color{codefg}
    }
}

% Configure hyperref colors consistently
\AtBeginDocument{
% Override pandoc's hidelinks setting with consistent options
\hypersetup{
    colorlinks=true,
    allcolors=red,
    linkcolor=red,
    urlcolor=red,
    citecolor=red,
    filecolor=red,
    menucolor=red,
    linktoc=all
}
}

% Simple page break support for document structure
\usepackage{bookmark}
\IfFileExists{xurl.sty}{\usepackage{xurl}}{} % add URL line breaks if available
\urlstyle{same}
% fallback for those not using the hyperref driver hyperxmp:
\makeatletter
\@ifundefined{xmpquote}{\newcommand{\xmpquote}[1]{#1}}{}
\makeatother
\hypersetup{
  colorlinks=true,
  linkcolor={red},
  filecolor={red},
  citecolor={red},
  urlcolor={red},
  pdfcreator={LaTeX via pandoc}}

\title{Conversions and Specification}
\author{Daniel Ari Friedman\\ ORCID: 0000-0001-6232-9096\\ Email: daniel@activeinference.institute\\ DOI: 10.5281/zenodo.16887791}
\date{October 07, 2026}

\begin{document}
\maketitle

{
\hypersetup{linkcolor=black}
\setcounter{tocdepth}{3}
\tableofcontents
}
\setstretch{1.0}
\section{Conversions \& Specification}\label{conversions-specification}

\subsection{Overview}\label{overview}

Sections \href{02_4d_namespaces.md}{2}, \href{03_quadray_methods.md}{3}
and \href{13_lattice_tooling.md}{13} establish the Quadray/IVM framework
and its computational tooling; this section documents the conversion
layer that joins its two coordinate namespaces ---
\passthrough{\lstinline!conversions.py!}, the module that maps between
Fuller.4D quadray coordinates and Coxeter.4D Cartesian coordinates.
Three concerns live here: the embedding itself and its projective fiber,
the Gram identities that make lattice distances exact, and an
\textbf{exact rational inversion} that answers ``which lattice point is
this point?'' with \passthrough{\lstinline!fractions.Fraction!}
arithmetic instead of floating-point rounding.

The full mathematical specification of this layer is the repository-root
file \passthrough{\lstinline!SPEC.md!}; the
\passthrough{\lstinline!tests/test\_spec\_examples.py!} pins its worked
examples numerically. This section is the manuscript projection of that
specification: same definitions, same error taxonomy, same guarantees,
prose instead of source. The analytical foundations it draws on are
\href{03_quadray_methods.md}{Quadray Methods} (normalization, the
\passthrough{\lstinline!Quadray!} class), the shell machinery of
\href{11_ivm_field_learning.md}{IVM Field Learning}, and the
nearest-site search of \href{13_lattice_tooling.md}{Lattice Tooling},
whose exact-distance identity this section re-derives from the
embedding's Gram matrix.

\subsection{The embedding (Fuller.4D to
Coxeter.4D)}\label{the-embedding-fuller.4d-to-coxeter.4d}

The default embedding maps the four quadray axes \((A, B, C, D)\) to the
vertices of a regular tetrahedron in \(\mathbb{R}^3\). With scale factor
\(s\) (uniform; scales all resulting coordinates):

\begin{equation}
\label{eq:conv-embedding}
M \;=\; s \times
\begin{pmatrix}
 1 & -1 & -1 &  1 \\
 1 &  1 & -1 & -1 \\
 1 & -1 &  1 & -1
\end{pmatrix},
\qquad \mathrm{xyz}(q) \;=\; M\,q .
\end{equation}

\passthrough{\lstinline!urner\_embedding(scale=1.0)!} builds exactly
this matrix (each row scaled by \passthrough{\lstinline!scale!}), and
\passthrough{\lstinline!quadray.DEFAULT\_EMBEDDING!} is the same
unscaled matrix stored as a tuple-of-rows; the two agree entry for entry
at scale 1. Every row of \(M\) sums to zero:

\begin{equation}
\label{eq:conv-fiber}
\textstyle\sum_{j} M_{ij} \;=\; 0 \quad \forall\, i,
\qquad \text{hence} \qquad
q \;\sim\; q + t\,\mathbf{1}, \qquad t \in \mathbb{Z},
\end{equation}

where \(\mathbf{1} = (1, 1, 1, 1)\). The vector \(\mathbf{1}\) spans the
kernel of \(M\), so the map is projective:
\passthrough{\lstinline!to\_xyz(q) == to\_xyz(q + t*(1, 1, 1, 1))!} for
every integer \(t\), and quadray normalization (see
\href{03_quadray_methods.md}{Quadray Methods}) simply selects the min-0
representative of that equivalence class. The forward map
\passthrough{\lstinline!quadray\_to\_xyz(q, M=None)!} delegates to
\passthrough{\lstinline!quadray.to\_xyz!}; passing
\passthrough{\lstinline!M=None!} selects
\passthrough{\lstinline!quadray.DEFAULT\_EMBEDDING!}, and an explicit
\passthrough{\lstinline!M!} is shape-validated against the shared
internal helper \passthrough{\lstinline!\_embedding\_rows!} (a
\passthrough{\lstinline!(3, 4)!} array is required,
\passthrough{\lstinline!ValueError!} otherwise) before the product is
taken. Existing callers that pass an explicit embedding --- for example
\passthrough{\lstinline!quadmath/scripts/quadray\_clouds.py!} --- are
unaffected.

\subsection{Gram identities and exact
distances}\label{gram-identities-and-exact-distances}

The embedding matrix satisfies two exact Gram identities. Row-wise, at
scale \(s\):

\begin{equation}
\label{eq:conv-gram-row}
M\,M^{\mathsf{T}} \;=\; 4\,s^{2}\,I_{3},
\end{equation}

and column-wise, with \(C_j\) the four columns of \(M\) and \(J_4\) the
all-ones matrix:

\begin{equation}
\label{eq:conv-gram-col}
M^{\mathsf{T}}M \;=\; 4\,I_{4} \;-\; J_{4},
\qquad \text{i.e.} \qquad
C_i \cdot C_j \;=\; 4\,\delta_{ij} \;-\; 1 \;\; (s = 1).
\end{equation}

\passthrough{\lstinline!embedding\_basis(M=None)!} returns the columns
of the validated embedding as a \passthrough{\lstinline!(4, 3)!} array
\(B\) whose \emph{rows} are the \(C_j\); at scale 1 the identity
\eqref{eq:conv-gram-col} reads \(B\,B^{\mathsf{T}} = 4\,I_4 - J_4\)
exactly (diagonal 3, off-diagonal \(-1\)), and it scales by \(s^2\) for
\passthrough{\lstinline!urner\_embedding(scale)!}. This is the identity
that \passthrough{\lstinline!lattice\_search.squared\_distance!} relies
on: for integer quadrays \(p, q\) with \(\delta = p - q\), the Euclidean
squared distance in embedded coordinates is the exact integer

\begin{equation}
\label{eq:conv-distance}
d^{2}(p, q) \;=\; \sum_{i} (M\delta)_i^{2}
\;=\; 4 \sum_{j} \delta_j^{2} \;-\; \Bigl(\sum_{j} \delta_j\Bigr)^{2},
\end{equation}

the same identity documented as Eq. \eqref{eq:lattice-distance-identity}
in \href{13_lattice_tooling.md}{Lattice Tooling}, where it powers
ranking and filtering; the test suite cross-validates it against float
embeddings and \passthrough{\lstinline!quadray.distance!}. Because the
identity is exact on integers, nearest-site decisions never suffer
floating-point boundary errors.

Executable check of the two identities:

\begin{lstlisting}[language=Python]
import numpy as np
from quadmath.lattice.conversions import urner_embedding, embedding_basis
from quadmath.lattice.lattice_search import squared_distance

M = urner_embedding()
B = embedding_basis()                              # (4, 3): rows are the columns of M
assert np.allclose(M @ M.T, 4.0 * np.eye(3))       # Eq. (eq:conv-gram-row)
assert np.allclose(B @ B.T, 4.0 * np.eye(4) - np.ones((4, 4)))
delta = np.array([2, 1, 1, 0])                     # shell-1 displacement
assert squared_distance(delta, np.zeros((1, 4), dtype=np.int64)) == 8   # Eq. (eq:conv-distance)
\end{lstlisting}

\subsection{Exact canonical inversion}\label{exact-canonical-inversion}

The forward map is projective but not injective: by
\eqref{eq:conv-fiber} the fiber of an embedded point is the full coset
\(\{\,q + t\,\mathbf{1} : t \in \mathbb{Z}\,\}\), and an inverse must
select one representative from it.

\begin{equation}
\label{eq:conv-canonical}
y \;=\; (x_1, x_2, x_3, 0), \qquad
\begin{pmatrix} C_0 & C_1 & C_2 \end{pmatrix}
\begin{pmatrix} x_1 \\ x_2 \\ x_3 \end{pmatrix} \;=\; \mathrm{xyz},
\qquad q^\ast \;=\; \mathrm{normalize}(y).
\end{equation}

\passthrough{\lstinline!xyz\_to\_quadray\_canonical(xyz, M=None)!}
resolves that ambiguity deterministically by exact rational arithmetic
(\passthrough{\lstinline!fractions.Fraction!}; a Python float enters as
its exact binary rational --- no rounding, no tolerance). The algorithm,
in prose: (1) convert \passthrough{\lstinline!xyz!} and
\passthrough{\lstinline!M!} to exact \passthrough{\lstinline!Fraction!}
entries and require every row of \passthrough{\lstinline!M!} to sum to
exactly zero, so that normalization preserves the image point
\eqref{eq:conv-fiber}; (2) because \(C_3 = -(C_0 + C_1 + C_2)\), rank
\(M = 3\) is equivalent to \(\det\,[\,C_0\; C_1\; C_2\,] \neq 0\) (a
\passthrough{\lstinline!ValueError!} otherwise); (3) solve the
\(3 \times 3\) system on columns 0--2 by Cramer's rule for the
particular preimage \(y = (x_1, x_2, x_3, 0)\) of
\eqref{eq:conv-canonical}; (4) the embedded point lies in the image of
the integer lattice \(\{\,M\,q : q \in \mathbb{Z}^{4}\,\}\) if and only
if \(x_1, x_2, x_3\) are all integers --- then every integer preimage is
\(y + t\,\mathbf{1}\), and the min-0 representative
\passthrough{\lstinline!Quadray(x1, x2, x3, 0).normalize()!} is the
unique deterministic tie-break on the fiber \eqref{eq:conv-fiber} (the
same rule as \passthrough{\lstinline!Quadray.normalize!}); a
non-integral preimage raises \passthrough{\lstinline!ValueError!} (the
point is not in the image).

The taxonomy is fail-closed. \passthrough{\lstinline!ValueError!}
covers: \passthrough{\lstinline!xyz!} of length other than 3;
\passthrough{\lstinline!M!} of the wrong shape; any row sum different
from zero; \(\det = 0\) (rank below 3); a non-integral preimage.
\passthrough{\lstinline!TypeError!} covers entries that are not
\passthrough{\lstinline!int!}/\passthrough{\lstinline!float!}/\passthrough{\lstinline!Fraction!}/\passthrough{\lstinline!numbers.Integral!}
--- note that \passthrough{\lstinline!numpy.float64!} (a
\passthrough{\lstinline!float!} subclass) and
\passthrough{\lstinline!numpy.int64!} (integral) are accepted, while
\passthrough{\lstinline!numpy.float32!} is rejected. There is no
floating-point fuzz anywhere: an off-lattice point raises rather than
snapping to a neighbor.

The float-valued inverse
\passthrough{\lstinline!quadray.quadray\_from\_xyz!} is a different,
deliberately fuzzier contract: it applies the pseudoinverse
\(M^{+} = M^{\mathsf{T}}(MM^{\mathsf{T}})^{-1}\) and rounds each
component half-up (\passthrough{\lstinline!floor(v + 0.5)!}, not
banker's rounding --- its docstring argument shows why half-up provably
stays in the \(\mathbf{1}\)-coset of the exact preimage, where
per-component round-half-even can leave it). For points on the lattice
it round-trips exactly; for general \(\mathbb{R}^3\) points it returns
the nearest lattice point in quadray coordinates (component-wise --- not
always nearest in embedded XYZ distance), which is a snapping operation,
not an inverse. The component-wise rule is the chosen contract; an
XYZ-nearest variant would be a separate function that searches
neighboring classes.
\passthrough{\lstinline!quadray\_roundtrip(q, M=None)!} pins the
round-trip contract

\begin{equation}
\label{eq:conv-roundtrip}
\mathrm{from\_xyz}\bigl(\mathrm{to\_xyz}(q)\bigr) \;=\; q ,
\end{equation}

and raises \passthrough{\lstinline!AssertionError!} explicitly (not a
bare \passthrough{\lstinline!assert!}, so the check survives
\passthrough{\lstinline!python -O!}) when it fails. The contract is
exact for normalized integer quadrays and, because
\passthrough{\lstinline!quadray\_roundtrip!} threads the \emph{same}
embedding through both legs, it is \textbf{scale-robust}: for
\(M = c\,M_0\) with any \(c \neq 0\) the pseudoinverse projection is
scale-independent,

\begin{equation}
\label{eq:conv-roundtrip-scale}
\mathrm{pinv}(cM)\,(cM\,q) \;=\; q \;-\; \frac{\sum_i q_i}{4}\,\mathbf{1},
\end{equation}

since \(\mathrm{pinv}(cM) = (1/c)\,\mathrm{pinv}(M)\) and the row space
of \(cM_0\) is that of \(M_0\) (the \((1/c)\cdot c\) cancels). A shell
site embedded at scale \(1/2\) --- the half-integer FCC picture of the
same sites --- round-trips through the same scaled embedding. The single
documented \passthrough{\lstinline!AssertionError!} case is an
\textbf{unnormalized} input: the inverse canonicalizes to the min-0
representative, e.g.~\passthrough{\lstinline!(2, 2, 2, 1)!} comes back
as \passthrough{\lstinline!(1, 1, 1, 0)!}. Embeddings whose rows do not
sum to zero lie outside the Urner family of \eqref{eq:conv-embedding}
and outside this guarantee.

Executable check of exact recovery and round-trips:

\begin{lstlisting}[language=Python]
from fractions import Fraction
from quadmath.lattice.conversions import (
    quadray_to_xyz, xyz_to_quadray_canonical, quadray_roundtrip, urner_embedding,
)
from quadmath.core.quadray import Quadray
from quadmath.lattice.omni_numbering import sites_through_shell

q0 = Quadray(2, 1, 1, 0)                          # cuboctahedron vertex, shell 1
assert xyz_to_quadray_canonical(quadray_to_xyz(q0)) == q0
assert xyz_to_quadray_canonical(
    (Fraction(-1), Fraction(1), Fraction(1))
) == Quadray(1, 1, 1, 0)                          # exact entries, no float fuzz
for site in sites_through_shell(2)[:8]:           # round-trips over enumerated sites
    assert quadray_roundtrip(Quadray(*site)) == Quadray(*site)
assert quadray_roundtrip(q0, urner_embedding(0.5)) == q0   # scale-robust (same rows, both legs)
\end{lstlisting}

\subsection{Lattice context}\label{lattice-context}

The canonical representatives produced by \eqref{eq:conv-canonical} sit
in the IVM site structure of \href{11_ivm_field_learning.md}{IVM Field
Learning}. A normalized quadray whose component sum is
\(\equiv 0 \pmod 4\) is an IVM sphere center
(\passthrough{\lstinline!is\_ivm\_site!}); the other three residue
classes are the octahedral and tetrahedral voids of the packing. The
radial observable is the shell norm

\begin{equation}
\label{eq:conv-shell}
N(q) \;=\; \sum_{i} \Bigl|\, q_i \;-\; \sigma/4 \,\Bigr| \;=\; 2k,
\qquad \sigma \;=\; \sum_{i} q_i ,
\end{equation}

computed by \passthrough{\lstinline!quadray\_shell\_norm!}: shell \(k\)
is the set of sites with \(N(q) = 2k\), populated by \(10k^2 + 2\)
centers (Eq. \eqref{eq:lattice-shell-population} in
\href{13_lattice_tooling.md}{Lattice Tooling}), and
\passthrough{\lstinline!ivm\_field.shell\_sites(k)!} enumerates shell
\(k\) in lexicographic order (1, 12, 42, 92, 162 sites for
\(k = 0..4\)). Conversions, shells, and nearest-site search share one
geometry; Sections \href{11_ivm_field_learning.md}{11} and
\href{13_lattice_tooling.md}{13} detail the field learner and the query
machinery respectively.

\subsection{API Summary}\label{api-summary}

{\def\LTcaptype{none} % do not increment counter
\begin{longtable}[]{@{}
  >{\raggedright\arraybackslash}p{(\linewidth - 2\tabcolsep) * \real{0.5000}}
  >{\raggedright\arraybackslash}p{(\linewidth - 2\tabcolsep) * \real{0.5000}}@{}}
\toprule\noalign{}
\begin{minipage}[b]{\linewidth}\raggedright
Function
\end{minipage} & \begin{minipage}[b]{\linewidth}\raggedright
Purpose
\end{minipage} \\
\midrule\noalign{}
\endhead
\bottomrule\noalign{}
\endlastfoot
\passthrough{\lstinline!conversions.urner\_embedding(scale=1.0)!} & The
\((3, 4)\) embedding matrix of Eq. \eqref{eq:conv-embedding} \\
\passthrough{\lstinline!conversions.quadray\_to\_xyz(q, M=None)!} &
Forward map via \passthrough{\lstinline!quadray.to\_xyz!};
\passthrough{\lstinline!M=None!} =
\passthrough{\lstinline!quadray.DEFAULT\_EMBEDDING!} \\
\passthrough{\lstinline!conversions.xyz\_to\_quadray\_canonical(xyz, M=None)!}
& Exact rational canonical inverse, Eqs. \eqref{eq:conv-fiber} and
\eqref{eq:conv-canonical} \\
\passthrough{\lstinline!conversions.quadray\_roundtrip(q, M=None)!} &
Asserts the round-trip identity \eqref{eq:conv-roundtrip}, scale-robust
per \eqref{eq:conv-roundtrip-scale} \\
\passthrough{\lstinline!conversions.embedding\_basis(M=None)!} &
\passthrough{\lstinline!(4, 3)!} array of embedding columns; Gram
identity \eqref{eq:conv-gram-col} \\
\passthrough{\lstinline!quadray.to\_xyz(q, embedding)!} & Underlying
forward product \(\mathrm{xyz} = M\,q\) \\
\passthrough{\lstinline!quadray.quadray\_from\_xyz(x, y, z, embedding)!}
& Float inverse: pseudoinverse + round-half-up (snapping, not exact
inversion) \\
\passthrough{\lstinline!lattice\_search.squared\_distance(p, sites)!} &
Exact squared distances via Eq. \eqref{eq:conv-distance} \\
\end{longtable}
}

\subsection{Verification}\label{verification}

\passthrough{\lstinline!tests/test\_conversions.py!} covers the layer
end to end: round-trips of \eqref{eq:conv-roundtrip} over shells 0--4
(309 enumerated sites, center included) and at scaled embeddings per
\eqref{eq:conv-roundtrip-scale}, exact canonical recovery of every site
through \eqref{eq:conv-canonical}, the Gram identities
\eqref{eq:conv-gram-row} and \eqref{eq:conv-gram-col} with
\passthrough{\lstinline!embedding\_basis!}, the distance identity
\eqref{eq:conv-distance} cross-checked against float embeddings and
\passthrough{\lstinline!quadray.distance!}, and the fail-closed error
taxonomy (\passthrough{\lstinline!ValueError!} for off-image points, bad
shapes, zero determinants and non-zero row sums;
\passthrough{\lstinline!TypeError!} for inexact-foreign entry types such
as \passthrough{\lstinline!numpy.float32!}). The repository-root
\passthrough{\lstinline!SPEC.md!} (the full mathematical specification)
and \passthrough{\lstinline!tests/test\_spec\_examples.py!} pin the
worked examples of this section numerically as they land. Run the suite:

\begin{lstlisting}[language=bash]
PYTEST_DISABLE_PLUGIN_AUTOLOAD=1 uv run coverage run -m pytest -q
uv run coverage report
\end{lstlisting}


\end{document}
